The transposition is an argument by identical premises, and its premises are exactly the imported and declared ones: indifference (Fact \claimref{one-relation.imported-indifference-instrument}, a premise), superselection-as-darkness (Fact \claimref{one-relation.imported-flag-arrow-superselection}), and irreversible recording (Premise \claimref{one-relation.recording-irreversible-stated-so-it-can}). The counting argument is a counting statement about the framework's own moduli: the trichotomy classifies the continuous moduli as two, on the stance recorded in Fact \claimref{one-relation.imported-flag-arrow-superselection}, and the CPT covariance argument makes the discrete freedom one bit, under its three premises, so any mechanism requiring one fresh uncaused bit per measurement event demands unboundedly many bits, exceeding the counted total; this is arithmetic on imported counts and is not scripted. The two gaps are stated, not proved away, which is the proposition's level: structural-conditional, resting on the dark-matter paper's asymmetry stance (the two numbers), on the open Hypothesis (S) (the one bit), and on this paper's own recording premise (the transposition).
